\documentclass[11pt,a4paper]{article}
\usepackage[T1]{fontenc}
\usepackage[utf8]{inputenc}
\usepackage[main=italian,provide=*]{babel}
\usepackage{amsmath,amssymb}
\usepackage{geometry}
\geometry{margin=2.3cm,top=2.5cm,bottom=2.7cm}
\usepackage{booktabs}
\usepackage[table,dvipsnames]{xcolor}
\usepackage{tcolorbox}
\tcbuselibrary{skins,breakable}
\usepackage{titlesec}
\usepackage{enumitem}
\usepackage{needspace}
\usepackage{fancyhdr}
\usepackage{graphicx}
\usepackage{hyperref}
\usepackage{longtable}
\hypersetup{colorlinks=true,linkcolor=Sepia,citecolor=Sepia,urlcolor=Sepia}
\newcommand{\Tr}{\operatorname{Tr}}

\definecolor{inkblue}{HTML}{1B3A5C}
\definecolor{lightblue}{HTML}{EAF1F8}
\definecolor{accent}{HTML}{B0562A}
\definecolor{softgray}{HTML}{6B6B6B}
\definecolor{okgreen}{HTML}{2E6B3E}
\definecolor{okgreenbg}{HTML}{EAF5EC}
\definecolor{warnamber}{HTML}{9C6B12}
\definecolor{warnbg}{HTML}{FBF1E0}

\titleformat{\section}
  {\normalfont\Large\bfseries\color{inkblue}}
  {\thesection}{0.6em}{}[{\color{inkblue!40}\titlerule}]
\titlespacing{\section}{0pt}{0.9em}{0.4em}
\titleformat{\subsection}
  {\normalfont\large\bfseries\color{accent}}
  {}{0em}{}
\titlespacing{\subsection}{0pt}{0.6em}{0.2em}

\pagestyle{fancy}
\fancyhf{}
\renewcommand{\headrulewidth}{0.4pt}
\fancyhead[L]{\small\color{softgray}Guida d'uso --- pacchetto riproducibile, rumore sul dimero}
\fancyhead[R]{\small\color{softgray}\thepage}
\fancyfoot[C]{\small\color{softgray}Tesi triennale, Samuele Schiavi --- Relatore: Prof.\ Alessandro Chiesa}

\newtcolorbox{keyeq}[1][]{
  colback=lightblue, colframe=inkblue!70, boxrule=0.6pt,
  arc=2pt, left=8pt, right=8pt, top=3pt, bottom=3pt, #1
}
\newtcolorbox{okbox}[1]{
  colback=okgreenbg, colframe=okgreen!60, boxrule=0.6pt,
  arc=2pt, left=7pt, right=7pt, top=2pt, bottom=2pt,
  fonttitle=\bfseries\color{okgreen}, title={#1}
}
\newtcolorbox{warnbox}[1]{
  colback=warnbg, colframe=warnamber!70, boxrule=0.6pt,
  arc=2pt, left=7pt, right=7pt, top=2pt, bottom=2pt,
  fonttitle=\bfseries\color{warnamber}, title={#1}
}

\newcommand{\code}[1]{\texttt{\small #1}}

\title{\vspace{-1.5em}\color{inkblue}Guida d'uso\\[0.2em]
\Large\normalfont\color{softgray}Riprodurre la pipeline di rumore sul dimero in locale}
\author{Pacchetto riproducibile --- Tesi triennale, Samuele Schiavi\\ Relatore: Prof.\ Alessandro Chiesa}
\date{}

\begin{document}
\sloppy
\maketitle
\thispagestyle{fancy}
\vspace{-2.0em}

Questo pacchetto contiene tutto il codice necessario per riprodurre, da
zero, ogni simulazione, ogni numero e ogni figura usati nei documenti
sulla Parte 2 (rumore quantistico) del dimero --- senza dipendere da
nessun file già calcolato: tutto si rigenera in locale. Ogni modulo e
ogni funzione del pacchetto sono commentati esplicitando sia il loro
ruolo all'interno del proprio file sia il loro ruolo nell'insieme della
pipeline --- questa guida ne è la mappa, non la ripetizione.

\tableofcontents

\section{Struttura del pacchetto}

\begin{keyeq}
\centering
\begin{tabular}{@{}ll@{}}
\code{codice/} & 16 moduli Python (mattoni + pipeline + estensioni)\\
\code{dati/} & dati intermedi (\code{.npz}) generati dagli script 01--03\\
\code{figure/} & le 15 figure finali, in PDF e PNG\\
\code{genera\_figure/} & 13 script, uno o più per ciascuna figura\\
\code{01\_genera\_dati\_vqe\_ideale.py} & primo script da eseguire\\
\code{02\_genera\_dati\_vqe\_noise\_aware.py} & secondo script\\
\code{03\_genera\_dati\_secondo\_punto.py} & terzo script (punto di lavoro alternativo)\\
\code{esegui\_tutto.py} & esegue tutto in un solo comando\\
\end{tabular}
\end{keyeq}

\subsection{I 16 moduli di \texttt{codice/}}

{\footnotesize
\begin{longtable}{@{}p{4.9cm}p{9.3cm}@{}}
\toprule
\textbf{File} & \textbf{Cosa fa} \\
\midrule
\endhead
\raggedright\code{dimer\_exact.py} & Hamiltoniana e diagonalizzazione esatta --- il ``mattone zero'', sorgente unica per tutto il pacchetto \\
\raggedright\code{vqe\_test2.py} & ansatz PMA-2q.3 e ottimizzazione VQE ideale (nessun rumore) \\
\raggedright\code{circuito\_correlazioni\_\allowbreak dimero.py} & circuito di Hadamard test per i correlatori, a rumore nullo (versione con supporto preparazione VQE reale) \\
\raggedright\code{noise\_model\_dimero.py} & costruzione del \code{NoiseModel} (canale depolarizzante + readout, calibrazione \code{ibm\_torino}) \\
\raggedright\code{vqe\_dm\_rumoroso\_\allowbreak dimero.py} & VQE con termine DM sotto rumore (riuso dei parametri ideali) \\
\raggedright\code{trotter\_rumoroso\_\allowbreak dimero.py} & quantum simulation (Trotter) sotto rumore \\
\raggedright\code{correlatori\_rumorosi\_\allowbreak dimero.py} & correlatori dinamici sotto rumore (readout simmetrico) \\
\raggedright\code{scan\_parametri\_rumore\_\allowbreak dimero.py} & scan di $N^*$ al variare di $\varepsilon_{1q},\varepsilon_{2q}$ \\
\raggedright\code{vqe\_noise\_aware\_\allowbreak dimero.py} & estensione: ottimizzazione ripetuta del VQE sotto rumore \\
\raggedright\code{verifica\_passo5\_\allowbreak noise\_aware.py} & estensione: verifica su tutta la griglia dello scan \\
\raggedright\code{correlatori\_readout\_\allowbreak asimmetrico.py} & estensione: errore di lettura asimmetrico \\
\raggedright\code{stile.py} & stile grafico condiviso da tutti gli script di figura \\
\raggedright\code{validate\_correlatori\_\allowbreak rumorosi\_dimero.py} & controllo obbligatorio: limite rumore nullo, conteggio gate, formula readout \\
\raggedright\code{validate\_scan\_parametri\_\allowbreak rumore\_dimero.py} & controllo obbligatorio: nessuna regressione, monotonia, invarianza da $p_\text{readout}$ \\
\raggedright\code{validate\_vqe\_noise\_\allowbreak aware\_dimero.py} & controllo obbligatorio: 8 verifiche sull'estensione VQE noise-aware \\
\raggedright\code{validate\_correlatori\_\allowbreak readout\_asimmetrico.py} & controllo obbligatorio: 4 verifiche sull'estensione readout asimmetrico \\
\bottomrule
\end{longtable}
}

\section{Prerequisiti}

Testato con questa combinazione esatta di versioni (altre versioni
recenti di Qiskit 2.x dovrebbero funzionare, ma non sono state
verificate):

\begin{keyeq}
\centering
\begin{tabular}{@{}ll@{}}
Python & 3.12.3\\
qiskit & 2.5.2\\
qiskit-aer & 0.17.2\\
numpy & 2.4.4\\
scipy & 1.17.1\\
matplotlib & 3.10.8\\
\end{tabular}
\end{keyeq}

Installazione consigliata (ambiente virtuale):
\begin{keyeq}
\centering
\begin{tabular}{@{}l@{}}
\code{python3 -m venv venv}\\
\code{source venv/bin/activate\ \ \ \# Windows: venv\textbackslash Scripts\textbackslash activate}\\
\code{pip install qiskit qiskit-aer numpy scipy matplotlib}\\
\end{tabular}
\end{keyeq}

\begin{okbox}{Nessuna installazione LaTeX richiesta}
Le figure sono generate direttamente in PDF/PNG da matplotlib, senza
passare da un motore di rendering matematico esterno
(\code{mathtext.fontset: cm} di matplotlib, autosufficiente).
\end{okbox}

\section{Il modo più semplice: un solo comando}

\begin{keyeq}
\centering
\code{cd pacchetto\_riproducibile}\\
\code{python3 esegui\_tutto.py}
\end{keyeq}

Esegue in ordine: generazione dati (3 script), validazioni (4 script),
tutte le 15 figure (13 script). Si ferma al primo errore, indicando
esattamente quale fase e comando ha fallito.

\begin{okbox}{Tempo stimato}
10--20 minuti. Le fasi più lente sono gli script di figura che ripetono
ottimizzazioni VQE a più valori di rumore.
\end{okbox}

Se tutto va a buon fine, l'ultima riga stampata è:
\begin{keyeq}
\centering
\code{TUTTO COMPLETATO. 15 figure in figure/, dati in dati/.}
\end{keyeq}

\section{Passo per passo, con la mappa esatta delle dipendenze}

Diciassette script in tutto (3 dati, 4 validazioni, 13 di figura).
Invece di un elenco lineare, questa sezione li organizza per
\textbf{onde di dipendenza reale} --- verificate leggendo ogni
\code{np.load}/\code{np.savez} nel codice, non assunte. Alcune
dipendenze sono meno ovvie del previsto: \code{03\_} non ha bisogno di
\code{01\_}/\code{02\_} (calcola tutto da zero per il secondo punto), e
la figura del secondo punto ha bisogno solo di \code{03\_}, non degli
altri due script dati.

\subsection{Onda 1 --- nessuna dipendenza}

\begin{keyeq}[title={01\_genera\_dati\_vqe\_ideale.py}]
\textbf{Ingresso}: nessuno.\\
\textbf{Uscita}: \code{dati/ground\_state\_test2.npz} --- \code{b,J,D},
\code{E0\_exact, psi0\_exact}, \code{vqe\_params} (3 parametri PMA-2q.3),
\code{fidelity, E\_vqe}.\\
\textbf{Chi lo consuma}: quasi tutto il resto del pacchetto (Onda 2) ---
è il prerequisito più condiviso. Fanno eccezione solo \code{03\_} e le
figure che dipendono esclusivamente da \code{03\_} (vedi sotto).\\
\textbf{Valore atteso}: energia $\approx-3.57321145$, fedeltà $\approx1.0000000000$.\\
\textbf{Comando} (dalla radice del pacchetto):\\
\code{python3 01\_genera\_dati\_vqe\_ideale.py}
\end{keyeq}

\begin{keyeq}[title={03\_genera\_dati\_secondo\_punto.py}]
\textbf{Ingresso}: nessuno --- NON legge \code{ground\_state\_test2.npz}
né \code{vqe\_noise\_aware\_result.npz}: ottimizza da zero sia il VQE
ideale sia quello noise-aware, al secondo punto ($b/J=-0.18$, $D/J=1$).\\
\textbf{Uscita}: \code{dati/ground\_state\_punto2.npz} (analogo del file
di \code{01\_} ma per il secondo punto) e
\code{dati/vqe\_noise\_aware\_punto2.npz} (analogo del file di \code{02\_}).\\
\textbf{Chi lo consuma}: solo \code{genera\_fig\_secondo\_punto.py}.\\
\textbf{Valori attesi}: $\Delta E\approx-2.56\times10^{-8}$,
$N^*$ Trotter $8\to8$, $N^*$ correlatore $3\to3$.\\
\textbf{Comando} (dalla radice del pacchetto):\\
\code{python3 03\_genera\_dati\_secondo\_punto.py}
\end{keyeq}

\begin{keyeq}[title={genera\_fig\_modello\_rumore.py}]
\textbf{Ingresso}: nessuno --- figura puramente concettuale (contrazione
del vettore di Bloch), nessun dato del progetto richiesto.\\
\textbf{Uscita}: \code{figure/fig\_depolarizzazione\_bloch.pdf}.\\
\textbf{Chi la consuma}: nessuno (punto di arrivo).\\
\textbf{Comando} (dalla radice del pacchetto):\\
\code{python3 genera\_figure/genera\_fig\_modello\_rumore.py}
\end{keyeq}

\subsection{Onda 2 --- dipende SOLO da \texorpdfstring{\code{01\_}}{01\_} (ground\_state\_test2.npz)}

\begin{keyeq}[title={02\_genera\_dati\_vqe\_noise\_aware.py}]
\textbf{Ingresso}: \code{dati/ground\_state\_test2.npz} (\code{vqe\_params, psi0\_exact}).\\
\textbf{Uscita}: \code{dati/vqe\_noise\_aware\_result.npz} ---
\code{x\_noise\_aware}, \code{E\_noise\_aware, F\_noise\_aware},
\code{E\_ideale\_su\_rumore, F\_ideale\_su\_rumore}, \code{cnot\_invarianza},
dati del confronto sul correlatore (\code{N\_grid\_corr},
\code{vals\_corr\_ideale/na}, \code{N\_star\_corr\_ideale/na}).\\
\textbf{Chi lo consuma}: le 4 figure dell'Onda 3 (uniche che hanno
bisogno anche di questo file, oltre a \code{01\_}).\\
\textbf{Valori attesi}: $\Delta E\approx-1.16\times10^{-7}$,
$\Delta F\approx+1.76\times10^{-8}$, $N^*$ del correlatore $5\to5$.\\
\textbf{Comando} (dalla radice del pacchetto, dopo \code{01\_}):\\
\code{python3 02\_genera\_dati\_vqe\_noise\_aware.py}
\end{keyeq}

\paragraph{I 4 controlli di validazione} Tutti e quattro leggono
\textbf{solo} \code{ground\_state\_test2.npz} (copiato in
\code{codice/}) --- nessuno tocca l'output di \code{02\_} o \code{03\_},
nemmeno \code{validate\_vqe\_noise\_aware\_dimero.py} (che pure valida
l'estensione noise-aware): ricalcola tutto da zero invece di leggere una
cache. Nessuno produce file, solo un resoconto a schermo
(\code{AssertionError} se un controllo fallisce); nessuno è consumato da
altri script (sono tutti controlli terminali). Prerequisito comune per
tutti e quattro, da fare una sola volta:\\
\code{cp dati/ground\_state\_test2.npz codice/}\\
\code{cd codice}

\begin{keyeq}[title={validate\_correlatori\_rumorosi\_dimero.py}]
\textbf{Comando}: \code{python3 validate\_correlatori\_rumorosi\_dimero.py}
\end{keyeq}

\begin{keyeq}[title={validate\_scan\_parametri\_rumore\_dimero.py}]
\textbf{Comando}: \code{python3 validate\_scan\_parametri\_rumore\_dimero.py}
\end{keyeq}

\begin{keyeq}[title={validate\_vqe\_noise\_aware\_dimero.py}]
\textbf{Comando}: \code{python3 validate\_vqe\_noise\_aware\_dimero.py}
\end{keyeq}

\begin{keyeq}[title={validate\_correlatori\_readout\_asimmetrico.py}]
\textbf{Comando}: \code{python3 validate\_correlatori\_readout\_asimmetrico.py}\\
Dopo l'ultimo controllo, torna alla radice: \code{cd ..}
\end{keyeq}

\paragraph{7 figure che dipendono solo da \texttt{01\_}} Ingresso comune
per tutte e sette: \code{dati/ground\_state\_test2.npz} soltanto. Uscita:
la rispettiva figura in \code{figure/}. Nessuna è consumata da altri
script del pacchetto (punti di arrivo). Comandi, tutti dalla
\textbf{radice} del pacchetto:

\begin{keyeq}[title={genera\_fig\_vqe\_dm\_rumoroso.py $\to$ fig\_passo2\_scan\_eps2q}]
\code{python3 genera\_figure/genera\_fig\_vqe\_dm\_rumoroso.py}
\end{keyeq}

\begin{keyeq}[title={genera\_fig\_trotter\_rumoroso.py $\to$ fig\_passo3\_F\_vs\_N, fig\_passo3\_zoom\_minimo}]
\code{python3 genera\_figure/genera\_fig\_trotter\_rumoroso.py}
\end{keyeq}

\begin{keyeq}[title={genera\_fig\_correlatori\_rumorosi.py $\to$ fig\_passo4\_correlatore\_vs\_N}]
\code{python3 genera\_figure/genera\_fig\_correlatori\_rumorosi.py}
\end{keyeq}

\begin{keyeq}[title={genera\_fig\_scan\_parametri.py $\to$ fig\_passo5\_Nstar\_vs\_eps}]
\code{python3 genera\_figure/genera\_fig\_scan\_parametri.py}
\end{keyeq}

\begin{keyeq}[title={genera\_fig\_trappola\_transpilazione.py $\to$ fig\_bug\_conteggio\_gate}]
\code{python3 genera\_figure/genera\_fig\_trappola\_transpilazione.py}
\end{keyeq}

\begin{keyeq}[title={genera\_fig\_scan\_eps2q\_confronto.py $\to$ fig\_scan\_eps2q\_confronto}]
Confronta ideale e noise-aware, ma ottimizza il noise-aware fresco ad
ogni valore di $\varepsilon_{2q}$ invece di leggere \code{02\_} --- per
questo resta in questa onda e non nella successiva.\\
\code{python3 genera\_figure/genera\_fig\_scan\_eps2q\_confronto.py}
\end{keyeq}

\begin{keyeq}[title={genera\_fig\_readout\_asimmetrico.py $\to$ fig\_correlatore\_N\_asimmetrico, fig\_Nstar\_vs\_rapporto}]
\code{python3 genera\_figure/genera\_fig\_readout\_asimmetrico.py}
\end{keyeq}

\subsection{Onda 2' --- dipende SOLO da \texorpdfstring{\code{03\_}}{03\_}}

\begin{keyeq}[title={genera\_fig\_secondo\_punto.py}]
\textbf{Ingresso}: \code{dati/ground\_state\_punto2.npz} e
\code{dati/vqe\_noise\_aware\_punto2.npz} --- NON tocca
\code{ground\_state\_test2.npz} né \code{vqe\_noise\_aware\_result.npz}:
è indipendente da \code{01\_} e \code{02\_}, dipende solo da \code{03\_}.\\
\textbf{Uscita}: \code{figure/fig\_secondo\_punto\_confronto.pdf}.\\
\textbf{Chi la consuma}: nessuno (punto di arrivo).\\
\textbf{Comando} (dalla radice del pacchetto, dopo \code{03\_}):\\
\code{python3 genera\_figure/genera\_fig\_secondo\_punto.py}
\end{keyeq}

\subsection{Onda 3 --- dipende da \texorpdfstring{\code{01\_}}{01\_} E \texorpdfstring{\code{02\_}}{02\_}}

\paragraph{4 figure che confrontano ideale e noise-aware} Ingresso
comune per tutte e quattro: \code{dati/ground\_state\_test2.npz}
\textbf{e} \code{dati/vqe\_noise\_aware\_result.npz} --- le uniche
figure di questo pacchetto che hanno davvero bisogno dell'output di
\code{02\_}. Uscita: la rispettiva figura in \code{figure/}. Nessuna
consumata da altri script. Comandi, dalla \textbf{radice} del pacchetto:

\begin{keyeq}[title={genera\_fig\_F\_vs\_N\_confronto.py $\to$ fig\_F\_vs\_N\_confronto}]
\code{python3 genera\_figure/genera\_fig\_F\_vs\_N\_confronto.py}
\end{keyeq}

\begin{keyeq}[title={genera\_fig\_correlatore\_confronto.py $\to$ fig\_correlatore\_noise\_aware\_confronto}]
\code{python3 genera\_figure/genera\_fig\_correlatore\_confronto.py}
\end{keyeq}

\begin{keyeq}[title={genera\_fig\_robustezza\_multistart.py $\to$ fig\_robustezza\_multistart}]
\code{python3 genera\_figure/genera\_fig\_robustezza\_multistart.py}
\end{keyeq}

\begin{keyeq}[title={genera\_fig\_griglia\_completa.py $\to$ fig\_passo5\_griglia\_completa}]
\code{python3 genera\_figure/genera\_fig\_griglia\_completa.py}
\end{keyeq}

\subsection{Il grafo delle dipendenze, in breve}

\begin{okbox}{Tre catene indipendenti, non una sola sequenza}
\begin{itemize}[leftmargin=1.3em,itemsep=0.15em]
\item \code{01\_} $\to$ (4 validate + 7 figure dirette) $\to$ \code{02\_} $\to$ (4 figure di confronto)
\item \code{03\_} $\to$ \code{genera\_fig\_secondo\_punto.py} (catena separata, mai tocca \code{01\_}/\code{02\_})
\item \code{genera\_fig\_modello\_rumore.py} (isolato, nessuna dipendenza)
\end{itemize}
\code{esegui\_tutto.py} le esegue comunque tutte in un ordine fisso
(\code{01\_}, \code{02\_}, \code{03\_}, poi validazioni, poi tutte le
figure in ordine alfabetico) --- più semplice da leggere nel log, anche
se non tutte le dipendenze lo richiederebbero.
\end{okbox}

\subsection{Tutti i comandi, in un unico blocco (nell'ordine di \texttt{esegui\_tutto.py})}

\begin{keyeq}
\centering
\begin{tabular}{@{}l@{}}
\code{cd pacchetto\_dimero\_rumoroso\_riproducibile}\\
\code{python3 01\_genera\_dati\_vqe\_ideale.py}\\
\code{python3 02\_genera\_dati\_vqe\_noise\_aware.py}\\
\code{python3 03\_genera\_dati\_secondo\_punto.py}\\
\code{cp dati/*.npz codice/}\\
\code{cd codice}\\
\code{python3 validate\_correlatori\_rumorosi\_dimero.py}\\
\code{python3 validate\_scan\_parametri\_rumore\_dimero.py}\\
\code{python3 validate\_vqe\_noise\_aware\_dimero.py}\\
\code{python3 validate\_correlatori\_readout\_asimmetrico.py}\\
\code{cd ..}\\
\code{python3 genera\_figure/genera\_fig\_F\_vs\_N\_confronto.py}\\
\code{python3 genera\_figure/genera\_fig\_correlatore\_confronto.py}\\
\code{python3 genera\_figure/genera\_fig\_correlatori\_rumorosi.py}\\
\code{python3 genera\_figure/genera\_fig\_griglia\_completa.py}\\
\code{python3 genera\_figure/genera\_fig\_modello\_rumore.py}\\
\code{python3 genera\_figure/genera\_fig\_readout\_asimmetrico.py}\\
\code{python3 genera\_figure/genera\_fig\_robustezza\_multistart.py}\\
\code{python3 genera\_figure/genera\_fig\_scan\_eps2q\_confronto.py}\\
\code{python3 genera\_figure/genera\_fig\_scan\_parametri.py}\\
\code{python3 genera\_figure/genera\_fig\_secondo\_punto.py}\\
\code{python3 genera\_figure/genera\_fig\_trappola\_transpilazione.py}\\
\code{python3 genera\_figure/genera\_fig\_trotter\_rumoroso.py}\\
\code{python3 genera\_figure/genera\_fig\_vqe\_dm\_rumoroso.py}\\
\end{tabular}
\end{keyeq}

Ciascuna validazione stampa un resoconto dettagliato con i valori
attesi accanto a quelli calcolati, e termina con un errore
(\code{AssertionError}) se un controllo fallisce --- non un semplice
messaggio ignorabile.


\section{Come sapere se un numero è quello giusto}

Ogni script stampa a schermo il valore appena calcolato \textbf{accanto}
al valore atteso (commento \code{\# atteso: ...} o riga di stampa
dedicata) --- non serve confrontare a memoria con i documenti.

\begin{warnbox}{Se un numero non coincide}
Qualcosa nel tuo ambiente locale (versione di Qiskit, seed del
generatore casuale, ordine di esecuzione) sta producendo un risultato
diverso: fermati e confronta prima di proseguire, non ignorare lo
scostamento.
\end{warnbox}

\section{Due insidie già trovate e corrette}

\begin{warnbox}{Non transpilare mai il circuito intero per i passi di Trotter ripetuti}
Se transpili l'intero circuito assemblato (ancilla + preparazione + $N$
passi + misura) a un livello di ottimizzazione alto, Qiskit può
riconoscere che gli $N$ passi identici formano un unico blocco e
comprimerli in un circuito a costo costante --- cancellando
silenziosamente la dipendenza da $N$ che è l'oggetto stesso dello
studio. La soluzione già in uso ovunque in questo pacchetto: transpilare
una volta il singolo passo, poi comporre il blocco già compilato $N$
volte senza più transpilare.
\end{warnbox}

\begin{warnbox}{Alcune funzioni hanno il punto di lavoro fissato come costante di modulo}
In particolare \code{vqe\_energia\_fedelta\_rumorosa()} (in
\code{vqe\_dm\_rumoroso\_dimero.py}) e \code{vqe\_noise\_aware()} (in
\code{vqe\_noise\_aware\_dimero.py}) costruiscono internamente
l'Hamiltoniana usando \code{b}, \code{D}, \code{J} definiti in testa al
modulo (il punto ``test 2''), non un argomento della funzione. Per
lavorare a un punto diverso (come nello script
\code{03\_genera\_dati\_secondo\_punto.py}) serve sovrascrivere
esplicitamente quegli attributi di modulo \textbf{prima} di chiamare la
funzione:
\begin{center}
\begin{tabular}{@{}l@{}}
\code{import vqe\_noise\_aware\_dimero as vna\_module}\\
\code{vna\_module.b, vna\_module.D, vna\_module.J = -0.18, 1.0, 1.0}\\
\end{tabular}
\end{center}
Altrimenti l'ottimizzazione avviene silenziosamente sul punto di lavoro
sbagliato, senza nessun errore --- solo numeri finali non corrispondenti
a quelli attesi. Questo pacchetto lo fa già correttamente in
\code{03\_genera\_dati\_secondo\_punto.py}; se scrivi un nuovo script per
un terzo punto di lavoro, applica lo stesso schema.
\end{warnbox}

\section{Cosa NON è incluso in questo pacchetto}

I documenti \code{.tex}/\code{.pdf} che usano queste figure (i cinque
documenti della pipeline base, l'estensione VQE noise-aware,
l'estensione readout asimmetrico) sono deliverable separati, già
consegnati in precedenza --- questo pacchetto riproduce i
\textbf{numeri e le figure}, non il testo che li commenta. I nomi dei
file delle figure qui prodotte coincidono esattamente con quelli
richiamati nei documenti (\code{\textbackslash includegraphics\{figure/nome.pdf\}}),
quindi puoi copiare la cartella \code{figure/} accanto a un documento
\code{.tex} per ricompilarlo.

\begin{okbox}{Verificato end-to-end, non solo assemblato}
L'intera pipeline è stata rieseguita da zero (cartelle \code{dati/} e
\code{figure/} svuotate, poi \code{esegui\_tutto.py} rilanciato) dopo
ogni modifica sostanziale al codice, incluse quelle di commento --- ogni
numero riportato in questa guida è stato ricontrollato, non solo
riportato dalla prima stesura.
\end{okbox}

\end{document}
